% ===== PRÉAMBULE CANONIQUE — à coller VERBATIM en tête de CHAQUE .tex =====
\documentclass[11pt,a4paper]{article}
\usepackage[utf8]{inputenc}\usepackage[T1]{fontenc}\usepackage{lmodern}
\usepackage[french]{babel}
\usepackage{amsmath,amssymb,mathtools}\usepackage{icomma}
\usepackage[version=4]{mhchem}          % \ce{...} : équations & formules
\usepackage{siunitx}\sisetup{locale=FR} % \qty{}{}, \si{}, \ang{}
\usepackage{chemfig}                    % \chemfig{...} : structures développées
\usepackage{xcolor}\usepackage{colortbl}
\usepackage{tikz}\usepackage{pgfplots}\pgfplotsset{compat=1.18}
\usepackage[most]{tcolorbox}\usepackage{enumitem}\usepackage{array,booktabs}
\usepackage[a4paper,margin=2.2cm]{geometry}
\usetikzlibrary{arrows.meta,calc,angles,quotes,positioning,patterns,backgrounds,shapes.geometric}
\definecolor{primary}{RGB}{222,49,99}\definecolor{primarydark}{RGB}{150,22,55}
\definecolor{defcol}{RGB}{0,114,178}\definecolor{methcol}{RGB}{52,92,148}
\definecolor{excol}{RGB}{0,135,90}\definecolor{attncol}{RGB}{200,40,55}
\tcbset{boxbase/.style={enhanced,breakable,boxrule=0.9pt,arc=2.5pt,
  left=3mm,right=3mm,top=2.3mm,bottom=2.3mm,fonttitle=\bfseries,coltitle=white}}
% --- boîtes TUTORIEL (noms EXACTS, ne pas renommer) ---
\newtcolorbox{cle}[1][]{boxbase,colback=defcol!6,colframe=defcol,
  title={\textsc{À retenir}\ifx\relax#1\relax\relax\else\textmd{ : \itshape #1}\fi}}
\newtcolorbox{methode}[1][]{boxbase,colback=methcol!7,colframe=methcol,sharp corners,
  title={\textsc{Méthode}\ifx\relax#1\relax\relax\else\textmd{ --- \itshape #1}\fi}}
\newtcolorbox{exemplebox}[1][]{boxbase,colback=excol!5,colframe=excol,
  title={\textsc{Exemple}\ifx\relax#1\relax\relax\else\textmd{ : \itshape #1}\fi}}
\newtcolorbox{piege}[1][]{boxbase,colback=attncol!6,colframe=attncol,
  title={\textsc{Piège}\ifx\relax#1\relax\relax\else\textmd{ : \itshape #1}\fi}}
\newtcolorbox{remarque}[1][]{boxbase,colback=black!4,colframe=black!45,
  title={\textsc{Remarque}\ifx\relax#1\relax\relax\else\textmd{ : \itshape #1}\fi}}
% --- boîtes EXERCICE / CORRIGÉ (noms EXACTS, auto-comptées) ---
\newtcolorbox[auto counter]{exo}[1][]{boxbase,colback=primary!4,colframe=primary,
  title={\textsc{Exercice}~\thetcbcounter\ifx\relax#1\relax\relax\else\textmd{ --- \itshape #1}\fi}}
\newtcolorbox[auto counter]{corrige}[1][]{boxbase,colback=excol!4,colframe=excol,
  title={\textsc{Corrigé}~\thetcbcounter\ifx\relax#1\relax\relax\else\textmd{ --- \itshape #1}\fi}}
\newcommand{\R}{\mathbb{R}}\newcommand{\N}{\mathbb{N}}\newcommand{\Z}{\mathbb{Z}}\newcommand{\ds}{\displaystyle}
% =========================================================================
\begin{document}
\sisetup{per-mode=symbol}
\begin{corrige}[pH de début de précipitation de \ce{Al(OH)3}]
La précipitation débute quand $K_{\mathrm{ps}}=[\ce{Al^{3+}}][\ce{OH-}]^{3}$. On isole $[\ce{OH-}]$ :
\[ [\ce{OH-}] = \sqrt[3]{\dfrac{K_{\mathrm{ps}}}{[\ce{Al^{3+}}]}}
             = \sqrt[3]{\dfrac{\num{2.0e-32}}{\num{1.0e-3}}} = \sqrt[3]{\num{2.0e-29}}
             \approx \qty{2.7e-10}{\mole\per\litre}. \]
\[ \mathrm{pH} = 14 + \log(\num{2.7e-10}) \approx 14 - 9{,}6 = 4{,}4. \]
\ce{Al(OH)3} précipite dès $\mathrm{pH}\approx 4{,}4$.
\end{corrige}

\begin{corrige}[pH de début de précipitation de \ce{Mg(OH)2}]
Type $1{:}2$ : $K_{\mathrm{ps}}=[\ce{Mg^{2+}}][\ce{OH-}]^{2}$, d'où
\[ [\ce{OH-}] = \sqrt{\dfrac{K_{\mathrm{ps}}}{[\ce{Mg^{2+}}]}}
             = \sqrt{\dfrac{\num{9.0e-12}}{0{,}010}} = \sqrt{\num{9.0e-10}} = \qty{3.0e-5}{\mole\per\litre}. \]
\[ \mathrm{pH} = 14 + \log(\num{3.0e-5}) \approx 14 - 4{,}5 = 9{,}5. \]
\ce{Mg(OH)2} ne précipite qu'en milieu franchement basique ($\mathrm{pH}\approx 9{,}5$).
\end{corrige}

\begin{corrige}[dissolution des sels d'acides faibles en milieu acide]
\begin{enumerate}[label=\alph*)]
  \item $\ce{CaCO3(s) + 2 H3O+ -> Ca^{2+}(aq) + CO2(g) + 3 H2O}$. L'ion \ce{CO3^{2-}} est consommé par
        \ce{H3O+} et transformé en \ce{CO2} gazeux qui s'échappe : l'équilibre de dissolution est
        déplacé vers la droite, la solubilité augmente.
  \item $\ce{FeS(s) + 2 H3O+ -> Fe^{2+}(aq) + H2S(g) + 2 H2O}$. L'ion \ce{S^{2-}} est consommé et
        part sous forme de \ce{H2S} gazeux : la dissolution est favorisée.
\end{enumerate}
\end{corrige}

\begin{corrige}[\ce{Fe(OH)3} précipite-t-il à pH $=3$ ?]
\begin{enumerate}[label=\alph*)]
  \item $[\ce{OH-}] = 10^{\,\mathrm{pH}-14} = 10^{3-14} = 10^{-11} = \qty{1.0e-11}{\mole\per\litre}$.
  \item $Q = [\ce{Fe^{3+}}][\ce{OH-}]^{3} = (\num{1.0e-2})(\num{1.0e-11})^{3}
           = (\num{1.0e-2})(\num{1.0e-33}) = \num{1.0e-35}$. Comme
        $Q=\num{1.0e-35} \gg K_{\mathrm{ps}}=\num{2.0e-39}$, \ce{Fe(OH)3} \textbf{précipite déjà à
        $\mathrm{pH}=3$} : c'est un hydroxyde extrêmement insoluble, il faut un milieu très acide pour
        le maintenir dissous.
\end{enumerate}
\end{corrige}

\begin{corrige}[solubilité d'un hydroxyde à pH imposé]
$s = K_{\mathrm{ps}}/[\ce{OH-}]^{2}$ avec $[\ce{OH-}]=10^{\mathrm{pH}-14}$.
\begin{enumerate}[label=\alph*)]
  \item $\mathrm{pH}=9$ : $[\ce{OH-}]=10^{-5}$, donc
        $s=\dfrac{\num{9.0e-12}}{(\num{1.0e-5})^{2}}=\dfrac{\num{9.0e-12}}{\num{1.0e-10}}=\qty{9.0e-2}{\mole\per\litre}$.
  \item $\mathrm{pH}=11$ : $[\ce{OH-}]=10^{-3}$, donc
        $s=\dfrac{\num{9.0e-12}}{(\num{1.0e-3})^{2}}=\dfrac{\num{9.0e-12}}{\num{1.0e-6}}=\qty{9.0e-6}{\mole\per\litre}$.
\end{enumerate}
\textbf{Constat :} en montant le pH de 9 à 11, la solubilité est divisée par \num{10000} — l'hydroxyde
précipite quand on alcalinise.
\end{corrige}
\end{document}
